% ECONOMICS_PROBLEM: {"id": "G21-262", "title": "Predicting and preventing bank runs", "jel_code": "G21", "source_id": "OP-0095"}
% !TeX program = xelatex
\documentclass[11pt,a4paper]{article}
\usepackage[margin=23mm]{geometry}
\usepackage{fontspec}
\setmainfont{Latin Modern Roman}
\usepackage{amsmath,amssymb,mathtools}
\usepackage{xcolor,enumitem,booktabs,longtable,array,xurl,hyperref}
\definecolor{muted}{HTML}{53636C}
\hypersetup{colorlinks=true,urlcolor=blue,linkcolor=blue}
\setlength{\parindent}{0pt}
\setlength{\parskip}{5pt}
\newcommand{\R}{\mathbb R}
\newcommand{\E}{\mathbb E}
\newcommand{\N}{\mathbb N}
\newcommand{\ind}{\mathbf 1}
\newcommand{\Var}{\operatorname{Var}}
\newcommand{\Cov}{\operatorname{Cov}}
\newcommand{\supp}{\operatorname{supp}}
\DeclareMathOperator*{\argmin}{arg\,min}
\DeclareMathOperator*{\argmax}{arg\,max}
\newcommand{\entrymeta}[3]{{\sffamily\small\color{muted}\textbf{\texttt{#1}}\quad|\quad #2\par #3\par}}
\newcommand{\Problemref}[1]{\texttt{#1}}
\newcommand{\JELref}[2]{\texttt{#2} (JEL \texttt{#1})}
\begin{document}
\section*{Predicting and preventing bank runs}
\textbf{Problem ID:} \texttt{G21-262}\quad \textbf{JEL:} \texttt{G21}\par
% BEGIN ECONOMICS STATEMENT
\label{problem:OP-0095}
\label{atlas:prob:F5}

\noindent\textbf{Atlas ID:} F5; \textbf{original number:} 95; \textbf{field:} Finance. \textbf{Classification:} Editorial research formulation (theoretical/empirical); not a certified unsolved theorem.

\subsection*{Definitions and assumptions}
Consider a bank issuing deposits of total face value $D>0$, with liquid reserves $L\ge0$ and an illiquid asset whose payoff depends on fundamental $z$. Depositors have privately observed liquidity needs and signals about $z$; they choose withdrawal or continuation under a stated payment-priority rule. Regulation $\pi$ specifies insurance, liquidity requirements and lender-of-last-resort access, including fiscal costs. A model $\theta$ gives the joint signal law, asset-sale prices and depositor network. An equilibrium selection determines withdrawals and repayment. Define a run event $A_H$ by a prespecified threshold of abnormal withdrawal or forced asset liquidation over horizon $H$.

\subsection*{Formal research question}
For publicly available information $\mathcal I_t$ and policy $\pi$, identify or bound
\[
 p_t(\theta,\pi)=\mathbb P_\theta(A_H\mid\mathcal I_t;\pi).
\]
Construct a prospective forecast that is calibrated over a stated bank population and evaluate its ability to predict systemic losses, not merely withdrawals. Characterize implementable policies minimizing expected real liquidation losses and fiscal resource costs subject to stated moral-hazard or coverage constraints. Determine which depositor-network, maturity and uninsured-funding measurements improve identification of equilibrium fragility before flight begins.

\subsection*{Interpretation and scope}
A bank-run probability does not follow from multiplicity alone: it requires a distribution over fundamentals, signals and an equilibrium-selection rule. A global-games perturbation can yield uniqueness under particular assumptions; it is not a general uniqueness theorem for all depositor networks or policy responses. Deposit withdrawals can be rational responses to insolvency or ordinary liquidity needs, so the event definition should distinguish these outcomes when relevant. Public forecasts may change beliefs and bank behavior; prospective evaluation must account for this intervention. Insurance reduces incentives to run but can alter ex ante risk taking and government exposure. Feasible policy evaluation includes asset quality, collateral haircuts and lending capacity.

\subsection*{Source audit and status}
[S1] establishes a classic liquidity-and-multiplicity model, [S2] supplies a selected-equilibrium approach under its assumptions, and [S3] provides depositor-level network evidence. Those benchmark results are not open. The present entry is an editorial prediction and policy-design extension; historical references do not certify its unresolved status in 2026.

\subsection*{References}

\begin{enumerate}[label={[S\arabic*]},leftmargin=*]

\item Douglas W. Diamond and Philip H. Dybvig (1983). \emph{Bank Runs, Deposit Insurance, and Liquidity}. Journal of Political Economy 91(3), 401–419. \url{https://doi.org/10.1086/261155}.
\textit{Locator and role:} Publisher or author abstract / bibliographic record; Liquidity provision and multiple equilibria in demand-deposit banking. Provides background or a benchmark result; does not establish that the editorial question remains unresolved in 2026.

\item Itay Goldstein and Ady Pauzner (2005). \emph{Demand-Deposit Contracts and the Probability of Bank Runs}. Journal of Finance 60(3), 1293–1327. \url{https://doi.org/10.1111/j.1540-6261.2005.00762.x}.
\textit{Locator and role:} Publisher or author abstract / bibliographic record; Global-games equilibrium selection and bank-run probability. Provides background or a benchmark result; does not establish that the editorial question remains unresolved in 2026.

\item Rajkamal Iyer and Manju Puri (2012). \emph{Understanding Bank Runs: The Importance of Depositor-Bank Relationships and Networks}. American Economic Review 102(4), 1414–1445. \url{https://doi.org/10.1257/aer.102.4.1414}.
\textit{Locator and role:} Publisher or author abstract / bibliographic record; Depositor-level withdrawal behavior, relationships and networks. Provides background or a benchmark result; does not establish that the editorial question remains unresolved in 2026.

\end{enumerate}

\subsection*{Common conventions: Source mathematical conventions}

Unless an entry states otherwise, agent and firm sets are finite. State and action spaces are measurable, strategies and outcome maps are measurable, probability laws are specified, and expectations exist for the targets under discussion. Infinite-horizon discount factors lie in $(0,1)$ and discounted objectives are finite. Norms, tolerances and complexity input sizes must be specified for computational comparisons. Constants or primitives that appear locally are inputs, not empirical facts asserted by this catalogue.

\textbf{Identification.} A statistical model $m\in\mathcal M$ implies an observable law $P_m$ and target $\theta(m)$. At law $P$, its identified set is
\[
\Theta(P;\mathcal M)=\{\theta(m):m\in\mathcal M,\ P_m=P\}.
\]
Point identification holds when this set is a singleton, or equivalently when $P_{m_1}=P_{m_2}$ implies $\theta(m_1)=\theta(m_2)$ for all compatible models. Sharp bounds mean bounds attained, or approached when the boundary is not attained, by compatible models. Writing this set is a definition, not a solution: the substantive task is to establish informative restrictions and derive or estimate the set. An unrestricted-confounding model may give only trivial bounds.

\textbf{Inference.} Identification, estimability and valid uncertainty quantification are separate questions. Fix $\alpha\in(0,1)$. A confidence procedure $C_n$ over a declared law class $\mathcal P$ has asymptotic uniform coverage at level $1-\alpha$ when
\[
\liminf_{n\to\infty}\inf_{P\in\mathcal P}\Pr_P\{\theta(P)\in C_n\}\geq1-\alpha.
\]
For set-identified targets the coverage event must be defined separately, for example coverage of the full identified set. If an entry uses identification-indexed models instead of uniquely indexed laws, the same coverage convention is applied over those models. Pointwise limits do not establish uniform validity, and asymptotic validity does not guarantee finite-sample precision. Sample sizes, dependence structures and weak-identification sequences are part of each application.

\textbf{Counterfactuals.} $Y_i(a)$ is a potential outcome under a specified intervention, including its timing, implementation and versions. It is not $Y_i$ conditional on voluntarily choosing $a$. A full intervention vector permits interference; shorthand scalar treatment notation does not silently impose no interference. Assignment, selection, attrition, anticipation and measurement must be specified. A claim about an instrument's local effect is not automatically a population policy effect.

\textbf{Economic equilibrium.} A counterfactual model includes best responses, constraints, feasibility, market clearing, state transitions and expectations consistent with the stated information. When several equilibria exist, either a declared selector is an input or the target is set-valued. Comparisons across policy regimes must use the stated selection convention. Reduced-form relationships are not presumed invariant after an equilibrium-changing intervention.

\textbf{Optimization and welfare.} Optimization is over the stated feasible set. If attainment is not established, $\sup$ or $\inf$ and $\varepsilon$-optimal choices replace an asserted maximizer or minimizer. For example, with value $V=\sup_{a\in\mathcal A}W(a)<\infty$, an $\varepsilon$-optimal set is $\{a\in\mathcal A:W(a)\geq V-\varepsilon\}$ for $\varepsilon>0$. Benefits, costs, risk and health or environmental outcomes can be aggregated only using declared units and conversion weights. Interpersonal utility weights, the treatment of fiscal transfers and foreign profits, discounting, and the behavioral welfare criterion are explicit inputs. A comparative-static derivative additionally requires existence and appropriate differentiability of the selected counterfactual.

\textbf{Sources.} Local labels [S1], [S2], and so forth restart in every question. Locators identify the relevant abstract, model, section or result at the level actually checked; a bibliographic or abstract-level verification is not represented as inspection of an entire proof. Working papers are distinguished from later publications and changing versions. Source summaries are paraphrased. All URL checks and editorial formulations refer to the audit date stated above.
% END ECONOMICS STATEMENT
\end{document}
